\documentclass[11pt]{article}
\usepackage[margin=1in]{geometry}
\usepackage{amsmath,amssymb,amsthm,mathrsfs}
\usepackage[T1]{fontenc}
\usepackage[colorlinks=true,linkcolor=blue,urlcolor=blue,citecolor=blue]{hyperref}
\newcommand{\Q}{\mathbf Q}
\newcommand{\Z}{\mathbf Z}
\newcommand{\F}{\mathbf F}
\newcommand{\K}{\mathscr K_E}
\newcommand{\Sha}{\operatorname{Sha}}
\newtheorem{theorem}{Theorem}
\newtheorem{lemma}{Lemma}
\newtheorem{corollary}{Corollary}
\title{An explicit reconstruction lattice for a\protect\\ N\'eron-normalized seed theta kernel}
\author{Dated supplement to the private BSD seed research release}
\date{24 September 2026}
\hypersetup{pdftitle={Neron-normalized seed theta kernel: reconstruction supplement},pdfauthor={BSD seed research project}}
\begin{document}
\maketitle
\begin{abstract}
For the fixed rank-two curve of the v0.2.2 dossier, we record a source-reviewed
written deduction that the N\'eron-normalized kernel of an integral degree-zero
theta packet lies in $(2^{10}\cdot4547)^{-1}\Z$. A comparison of congruence
ideals shows that its scale relative to the primitive quaternionic eigenvector
is a unit at every good prime $p\ge5$. These statements give a finite integer
acceptance test for the retained good-prime witness without computing the
unknown quaternionic norm. No seed kernel value or new primary component is
computed. The imported theorems and retained geometric inputs are explicit;
this supplement has model reviews, no independent human review and no new
Lean verification. It does not change the v0.2.2 manuscript or its claim grades.
\end{abstract}

\section{Objects and fixed normalizations}
The curve and conductor are
\[
 \begin{gathered}
 E:y^2=x^3-22735x^2+106272484x,\\
 N=2^3\cdot3\cdot2243\cdot4547^2\cdot5843.
 \end{gathered}
\]
The translation $x=X+7578$ has $u=1$ and gives the retained global minimal
model
\[
 [0,-1,0,-66020924,-65074271436].
\]
Thus $\omega_E=dx/(2y)$ is the
N\'eron differential. Existing claims A.1, A.3 and A.4 supply the local data,
full odd residual images and the chosen-map Manin bound respectively.

Let $N^-=3\cdot2243\cdot5843$ and $N^+=2^3\cdot4547^2$. In the definite
quaternion algebra ramified at $N^-$ and infinity, use the retained Eichler
order of level $N^+$ and its ideal-class set $X$. The stabilizer weights are
all one. Let $\psi\in\Z[X]$ be the primitive vector on the one-dimensional
$f_E$-line and put $M=\sum_{x\in X}\psi(x)^2$. For $z=\sum z_x[x]$ write
$\psi(z)=\sum z_x\psi(x)$ and $\deg z=\sum z_x$.

The retained theta convention is
\[
 \Theta_{x,y}(\tau)=\frac12\sum_{b\in L(x,y)}e^{2\pi iQ(b)\tau},
 \quad\Theta(z,w)=\sum_{x,y}z_xw_y\Theta_{x,y}.
\]
Here $L(x,y)$ is the normalized Hom lattice, even of rank four, with Gram
determinant $N^2$ and level $N$. With $B(1)=\mathrm{Id}$, the audited
prime-to-level Hecke projection satisfies
\begin{equation}\label{oldkernel}
 \lambda_f(\Theta(z,w))=\frac{\psi(z)\psi(w)}{M}.
\end{equation}
This is the retained rank-one theta identity, not a global identification
of $M$ with a classical congruence number or modular degree.

Let $E_*$ be the optimal quotient, $\phi_0:X_0(N)\to E_*$ its
parametrization of degree $d$, and $h:E_*\to E$ the identity or the
retained rational 2-isogeny. Set $e=\deg h\in\{1,2\}$ and
$\phi_E=h\phi_0$, with
\[
 \phi_E^*\omega_E=c_E\omega_f,\qquad
 \omega_f=2\pi i f(\tau)\,d\tau,\qquad c_E\mid2^9.
\]
The Manin bound refers to this chosen map, without further multiplication
maps. Its retained geometric dependencies are not reproved here.
Define
\[
 A_E=\frac i2\int_{E(\mathbf C)}\omega_E\wedge\overline{\omega_E},
 \qquad P_f=\int_{\Gamma_0(N)\backslash\mathfrak H}|f|^2\,dx\,dy.
\]
$A_E$ is the complex period-lattice area, not the BSD real period or the
N\'eron--Tate regulator. Pullback of area gives
\begin{equation}\label{area}
 ed\,A_E=4\pi^2c_E^2P_f,\qquad
 \K(G):=\frac{4\pi^2}{A_E}\langle G,f\rangle
       =\frac{ed}{c_E^2}\lambda_f(G).
\end{equation}
Consequently $\K(z,w):=\K(\Theta(z,w))=u_E\psi(z)\psi(w)$, where
$u_E=ed/(c_E^2M)$.

\section{The good-prime unit comparison (NK.1)}
\begin{theorem}[NK.1; written deduction with published imports]\label{unit}
For every prime $p\ge5$ with $p\nmid N$,
\[
 v_p(M)=v_p(d),\qquad u_E\in\Z_p^\times.
\]
In particular, for $z\in\Z[X]$, $\K(z,z)$ is a $p$-adic unit if and only
if $\psi(z)\not\equiv0\pmod p$.
\end{theorem}
\begin{proof}
We spell out the comparison chain, with the source hypotheses retained.
In the weight-two pairing of Kim--Ota \cite[\S6.1, (6.1)]{ko}, the
Atkin--Lehner element permutes $X$ and acts on its $f$-line by a sign.
All stabilizer weights are one. The pairing matrix on $\Z_p^X$ is therefore
a permutation matrix and is perfect. The globally primitive vector is
locally primitive, so the image of the pairing on its eigenline is
$M\Z_p$ (the sign is a unit). This identifies the quaternionic pairing
ideal $\xi_{f_B}$, not yet a classical congruence ideal.

The seed satisfies Kim--Ota's CR+ hypotheses. The weight is two and
$p\ge5$ is prime to $N$. The residual image is $\mathrm{GL}_2(\F_p)$;
its restriction to $\Q(\sqrt{p^*})$ contains $\mathrm{SL}_2(\F_p)$ and is
absolutely irreducible. The integer $N^-$ is squarefree with three prime
factors, and no prime exactly divides $N^+$. At $3,2243,5843$ the minimal
discriminant valuations $2,1,2$ are prime to $p$, so the residual Tate
inertia remains nontrivial. At 2 and 4547 inertia is finite of order
dividing 24 and of order 4 respectively. Its averaging idempotents are
$p$-integral, so invariant dimensions and Artin conductors are unchanged
on reduction. Thus $N(\bar\rho)=N$, and all the Tamagawa exponents in
\cite[Corollary 6.7]{ko} vanish. That corollary gives
\[
 v_p(M)=v_p(\xi_{f_B})
       =v_p(\eta_f(N^+,N^-))=v_p(\eta_f(N)).
\]
At nonordinary primes the freeness used here explicitly imports
\cite[Remark 6.6(1)]{ko}, the low-weight crystalline adaptation. Its
deformation-theoretic justification is not reproved in this supplement.

For completeness, the last classical lattice identification uses the
\emph{full} Hecke algebra $\mathbb T$ and the standard duality
$S_2(\Gamma_0(N),\Z)=\operatorname{Hom}_{\Z}(\mathbb T,\Z)$.
If $e_f$ is the rational idempotent, write
$\mathbb T\cap\Q e_f=\Z(ne_f)$ with $ne_f$ primitive. Then the
annihilator definition gives $\eta_f(N)=n\Z$ locally, while duality
shows that the projection of integral cuspforms is $n^{-1}\Z f$.
The quotient by $\Z f$ and its orthogonal complement is cyclic of order
$n$. Hence $v_p(\eta_f(N))=v_p(r)$ for the classical congruence number
$r$. The full-Hecke comparison and the localization are detailed in the
retained Phase 1 proof; they are not inferred solely from
\eqref{oldkernel}.

Agashe--Ribet--Stein \cite[Theorem 2.1]{ars} give $v_p(r)=v_p(d)$
for the optimal elliptic quotient at good $p$. Finally $e\in\{1,2\}$
and $c_E\mid2^9$ imply the asserted unit property of $u_E$.
\end{proof}
The conclusion is equality of valuations at the stated primes. It
asserts neither $p\nmid M$ nor a global equality $M=r$ or $M=d$.

\section{The explicit real reconstruction lattice (NK.2)}
\begin{lemma}\label{cusps}
For $z,w\in\Z[X]$ with $\deg z=0$, $G=\Theta(z,w)$ is cuspidal at every
cusp and integral at infinity. At a cusp denominator $c\mid N$, set
$n=v_q(N)$ and $a=v_q(c)$. Every positive Fourier coefficient in the
plain weight-two slash expansion has valuation at least $a-n$.
\end{lemma}
\begin{proof}
We use the retained Poisson expansion (lab Lemmas 4 and 4$'$). For
$\gamma=\left(\begin{smallmatrix}A&B\\c&D\end{smallmatrix}\right)$,
$c\ge1$, write $\mathrm e(t)=e^{2\pi it}$. The expansion is
\[
 (\Theta_{x,y}|_2\gamma)(\tau)
 =-\frac1{2Nc^2}\sum_{\ell\in L(x,y)^\#}
 S(A,\ell;c)\,\mathrm e(DQ(\ell)/c)\,\mathrm e(Q(\ell)\tau),
\]
\[
 S(A,\ell;c)=\sum_{r\in L(x,y)/cL(x,y)}
       \mathrm e\bigl((AQ(r)+(\ell,r))/c\bigr).
\]
There is no additional cusp-width factor in the slash operator.
The zero-vector Gauss value is independent of $(x,y)$, including its
local norm-unit choices. Every constant term of $G$ is therefore a
common scalar times $\deg z\deg w$ and vanishes. At infinity the
positive coefficients of each $\Theta_{x,y}$ are integers by pairing
$b$ and $-b$.

At an Eichler prime the two off-diagonal summations force their support
conditions and contribute $q^{2a}$. The remaining hyperbolic sum
contributes $q^a$ times a root of unity. Thus the local Gauss value is
zero or $q^{3a}$ times a root of unity. At a division prime, $n=1$;
the retained norm Gauss identity gives $-q^3$ times a root of unity on
support when $a=1$, and the factor is 1 when $a=0$. Pairing $\ell$
and $-\ell$ cancels the theta factor $1/2$ exactly, including at $q=2$.
The prefactor has valuation $-n-2a$, and other local factors are
$q$-units or zero. Hence each surviving summand, and its finite sum
at a positive Fourier frequency, has valuation at least $a-n$.
This is a lower bound, not a statement of integrality at every cusp.
\end{proof}

\begin{theorem}[NK.2; written deduction with retained Manin input]\label{lattice}
For $z,w\in\Z[X]$, $\deg z=0$, put $G=\Theta(z,w)$. Then
\[
 D_0\omega_G\text{ lies in the N\'eron cotangent lattice of }J_0(N),
 \qquad D_0=2\cdot4547,
\]
and
\[
 \boxed{B_\theta\K(G)\in\Z,\qquad
 B_\theta=2^{10}\cdot4547=4\,656\,128.}
\]
The weaker multiplier $B_{\mathrm{join}}=2^{17}\cdot4547=595\,984\,384$
also works.
\end{theorem}
\begin{proof}
CNS \cite[Proposition 5.14]{cns}, which applies to general cusp forms
in its stated domain, requires the respective valuation bounds
\[
 -n\quad(a=0),\qquad
 a-n+\frac1{q-1}\quad(0<a<n),\qquad0\quad(a=n).
\]
Lemma~\ref{cusps} meets these bounds at $q\in\{3,2243,5843\}$,
where $n=1$. At 2 ($n=3$) and 4547 ($n=2$), multiplying once by the
local prime meets all intermediate bounds. At good primes integrality
at infinity suffices. Thus $D_0\omega_G$ is a regular differential.
These are sufficient multiplier exponents; their optimality for the
actual packet forms is not established.

CNS \cite[Theorem 6.12]{cns} gives the required rational singularities
at every prime: at all primes at least five; at three since $v_3(N)=1$;
and at two since $3\mid N$ and $3\equiv3\pmod4$.
Proposition 6.2 and Example 6.3 of that paper identify regular
differentials with the N\'eron cotangent lattice. Local integrality
at all primes gives the asserted global lattice membership.

Let $\pi_E:J_0(N)\to E$ be the chosen quotient and $i_E:E\to J_0(N)$
its dual, so $\pi_Ei_E=[ed]$. Hecke equivariance kills the complement
of the $f$-line under $i_E^*$, while
$i_E^*\pi_E^*\omega_E=ed\omega_E$. Therefore
\[
 i_E^*\omega_G=\frac{ed}{c_E}\lambda_f(G)\omega_E
             =c_E\K(G)\omega_E.
\]
The map extends to N\'eron models, and $\omega_E$ generates the
integral cotangent lattice of $E$. It follows that
$D_0c_E\K(G)\in\Z$. The retained $c_E\mid512$ gives $B_\theta$.
The weaker multiplier follows also by using the retained whole-lattice
inclusion $2^8S\subseteq H$ at two. This latter argument avoids the
theta Gauss estimate at two; both arguments retain the same Manin bound.
\end{proof}

\section{A certificate specification, not a computed witness (NK.3)}
For an admissible good prime $p\ge5$, $p\nmid7N$, use the retained BLV
conductor-$p^2$ Gross packet. If $\gamma$ generates its first cyclic
layer, its mod-$p$ weight is
\[
 \overline z_p=\sum_{j=0}^{p-1}j^2[\gamma^jP_1],\qquad
 M_2(p)=\bar\psi(\overline z_p).
\]
Here the Gross point $P_1$ is not the Mordell--Weil basis point bearing
that name in the main paper. Equal trace degrees and
$\sum_{j\in\F_p}j^2=0$ make the mod-$p$ degree zero.
An integral lift $z_p$ need not have degree zero: balanced square
weights sum to 22 at $p=11$. Choose a degree-one class $[a]$ and set
\begin{equation}\label{lift}
 z_p^0=z_p-\deg(z_p)[a].
\end{equation}
Then $\deg z_p^0=0$ and $z_p^0\equiv\overline z_p\pmod p$.
The integral weights and anchor must be specified in a certificate.

\begin{corollary}[NK.3; conditional sufficient criterion]
For this lift define $n_p=B_\theta\K(z_p^0,z_p^0)$. Then
\[
 n_p\in\Z,\qquad
 p\nmid n_p\ \Longleftrightarrow\ M_2(p)\ne0
 \ \Longrightarrow\ \Sha(E/\Q)[p^\infty]=0.
\]
The last implication imports the existing ordinary/supersingular
forward bridge B.5; its antecedent is uncomputed here.
\end{corollary}
\begin{proof}
Theorem~\ref{lattice} gives the integer, and $p\nmid B_\theta$.
By Theorem~\ref{unit}, $n_p=B_\theta u_E\psi(z_p^0)^2$ has unit scale.
Equation~\eqref{lift} identifies its nonvanishing test with $M_2(p)$.
Apply B.5. No converse from vanishing of $\Sha$ to the witness is used.
\end{proof}

Let $[L,U]$ be a genuinely certified real interval for the \emph{final}
normalized scalar, including all numerical and truncation errors. Unique
integer recovery is possible when its radius is strictly less than
\[
 \frac1{2B_\theta}=\frac1{9\,312\,256};\qquad
 \text{fallback: }\frac1{2B_{\mathrm{join}}}=\frac1{1\,191\,968\,768}.
\]
The bound guarantees at most one lattice point; the theorem and enclosure
supply existence. For a wider interval put
$k_-=\lceil B_\theta L\rceil$ and $k_+=\lfloor B_\theta U\rfloor$.
A sufficient acceptance test is
\[
 k_-\le k_+,\qquad
 \left\lceil\frac{k_-}{p}\right\rceil>
 \left\lfloor\frac{k_+}{p}\right\rfloor.
\]
The first inequality rejects an empty lattice interval. The second excludes
every multiple of $p$. A nonunit or unresolved interval gives no conclusion
about $\Sha$. The sharp lattice cannot be applied to a noncuspidal anchor
diagonal merely because an anchored expression eventually produces this
cuspidal scalar.

\section{Evidence grade and remaining obligations}
The detailed proof sources, source hashes and three model-review records
are included with this supplement. The portable checks replay 262 exact
Eichler Gauss identities, 13 cusp-threshold comparisons and degree-zero
relifting, and compare 41,925 interval decisions with direct integer
enumeration. They verify finite algebra and acceptance arithmetic; they
do not formalize the source papers or evaluate a seed numerator.
The general division-place Gauss identity, the chosen-map Manin bound,
and Kim--Ota's nonordinary adaptation retain their stated written/import
grades. Numerical period checks are evidence only. There is no claim of
independent human review or of novelty relative to the literature.

The reconstruction prerequisite is now supplied for this normalized
degree-zero route. An affordable certified seed evaluation and a nonzero
seed witness are still missing. The fixed denominator $B_\theta$ is not
a fixed numerator controlling the prime support of $\Sha$: $n_p$ varies
with $p$ and its packet. No new primary component, FIN, ORDER9 or Q36
follows. Full BSD still requires an all-prime arithmetic conclusion and
the exact real leading-term comparison.

\begin{thebibliography}{9}
\bibitem{ko} C.-H. Kim and K. Ota,
\emph{On the quantitative variation of congruence ideals and integral
periods of modular forms}, source version
\href{https://arxiv.org/pdf/1905.02926v4}{arXiv:1905.02926v4},
\S6.1, Assumption 6.3, Theorem 6.5, Remark 6.6(1), Corollary 6.7.
\bibitem{ars} A. Agashe, K. Ribet and W. Stein,
\emph{The modular degree, congruence primes, and multiplicity one},
\href{https://www.math.fsu.edu/~agashe/math/moddeg3.pdf}{author-hosted manuscript},
Theorems 2.1 and 3.6, Remark 3.7.
\bibitem{cns} K. \v{C}esnavi\v{c}ius, M. Neururer and A. Saha,
\emph{The Manin constant and the modular degree}, source version
\href{https://arxiv.org/pdf/1911.09446v3}{arXiv:1911.09446v3},
Propositions 5.14 and 6.2, Example 6.3, Theorem 6.12.
\bibitem{lab} BSD seed lab, source commit
\href{https://github.com/Bome2017/bsd-seed-lab/commit/b358f7507759dad4a1e016c519d6f486ffca37e1}{\texttt{b358f7507759dad4a1e016c519d6f486ffca37e1}},
\emph{01\_NORMALIZATION\_REVIEW.md} and
\emph{02\_DENOMINATOR\_LATTICE.md}, 24 September 2026.
Copied with their context, source bindings and scoped reviews in this supplement.
\end{thebibliography}
\end{document}
